\chapter{Prototipo integrado}
\label{cap:8}

Este capítulo presenta de manera detallada el desarrollo e implementación del prototipo software concebido para validar la arquitectura de tres capas. A través de los componentes descritos, se muestra la traducción práctica de los modelos predictivos y las interfaces de datos en una aplicación funcional. De este modo, se demuestra la viabilidad técnica del sistema inteligente en un entorno local y reproducible.

\section{Introducción}

Este capítulo describe cómo se materializa en software la arquitectura definida en el Capítulo~\ref{cap:6}. El objetivo no es repetir el diseño conceptual, sino mostrar que el prototipo ejecuta una cadena completa: recibe o carga incidentes, estructura variables, aplica la Capa 2 de priorización, genera una explicación mediante la Capa 3 y registra trazabilidad para revisión posterior. Para ello, se examina el flujo de control desde que entra un suceso hasta que el operador valida la recomendación. Con esta descripción se busca cerrar la brecha entre los planteamientos conceptuales previos y la realidad del código fuente implementado.

El prototipo tiene alcance académico. No sustituye al sistema institucional del 112, no moviliza recursos, no activa planes y no toma decisiones autónomas. Su valor reside en demostrar que la propuesta puede ejecutarse de forma reproducible sobre datos reales abiertos y que sus salidas pueden auditarse.

\section{Arquitectura software}

El sistema se organiza en cinco bloques principales:

\begin{enumerate}
    \item \textbf{Pipeline de datos y Capa 1.} Limpia el Registro 112 Castilla y León, normaliza campos, construye variables y, en ejecución, extrae señales textuales predecisionales desde el texto operativo del incidente.
    \item \textbf{Capa 2 RuleFit.} Implementa supervisión débil, línea base experta, RuleFit-lite, comparativa con \texttt{imodels}, evaluación e inferencia.
    \item \textbf{Capa 3 explicativa.} Genera explicaciones a partir de la salida de la Capa 2, evidencias disponibles y corpus normativo recuperado mediante RAG/MCP. En la versión integrada se ejecuta localmente mediante Ollama.
    \item \textbf{Backend e interfaz de contratos.} Expone la inferencia y validación del flujo de IA mediante esquemas de validación FastAPI \cite{FastAPI2026} y una interfaz Streamlit \cite{Streamlit2026} de control del operador.
    \item \textbf{Trazabilidad y validación.} Registra salidas, métricas, muestra de revisión humana y artefactos de evaluación.
\end{enumerate}

La separación entre capas evita que el LLM decida la prioridad. La prioridad procede de la Capa 2; la Capa 3 traduce esa salida en una explicación revisable. Al desacoplar la lógica de predicción del procesamiento de lenguaje natural, se minimiza el impacto de posibles alucinaciones. Esta decisión arquitectónica refuerza la robustez del sistema frente a comportamientos no deterministas de la inteligencia artificial generativa.

\section{Componentes de Capa 2}

La parte de priorización se estructura en el módulo del motor de reglas y estimadores interpretables. Su lógica contiene la generación de etiquetas académicas P1--P4, la implementación de las reglas expertas que sirven como referencia y el cálculo comparativo con alternativas.

Esta estructura mantiene separado el modelo RuleFit-lite de la preparación de datos y de la lógica de explicación. También facilita comparar alternativas sin alterar el contrato de salida que recibe la Capa 3 y permite revisar cada componente de forma independiente.

En el prototipo integrado, el modelo seleccionado se empaqueta como un artefacto persistente y se carga durante la inferencia. Esta decisión permite separar el entrenamiento de la ejecución del demostrador y evita reentrenamientos accidentales. Además, el módulo incorpora comprobaciones de compatibilidad para asegurar la carga correcta de los parámetros del modelo, protegiendo la reproducibilidad del prototipo ante diferencias menores en el entorno de ejecución.

\section{Flujo de procesamiento}

El flujo funcional del prototipo es el siguiente:

\begin{enumerate}
    \item Se carga un incidente desde el dataset limpio o desde una entrada controlada.
    \item Se extraen variables operativas y señales textuales permitidas.
    \item Se verifica que no se usen columnas posteriores a la decisión (leakage).
    \item La Capa 2 calcula la prioridad P1--P4, probabilidades y reglas o señales activadas.
    \item La Capa 3 genera una explicación breve, apoyada en evidencias y corpus normativo cuando procede.
    \item El sistema registra la inferencia para auditoría y revisión humana.
\end{enumerate}

El resultado no es una orden operativa, sino una recomendación estructurada. El operador mantiene la capacidad de aceptar, modificar o rechazar la salida. Esta decisión de diseño preserva el principio de control humano en el lazo de decisión. La posible reducción del tiempo de respuesta o de la carga cognitiva deberá medirse en una validación posterior con personal operativo; no se considera demostrada por las pruebas actuales.

\section{Prototipo ejecutable}

El prototipo dispone de una ejecución integrada del flujo completo: texto del incidente, Capa 1 simbólica, Capa 2 RuleFit-lite, Capa 3 con LLM local y registro de la respuesta. El servicio de inferencia expone contratos de datos para comprobar la entrada y la salida, mientras que la interfaz permite introducir incidentes sintéticos por texto o por voz y visualizar la recomendación de Capa 2 junto con la explicación de Capa 3. El Anexo~\ref{cap:anexog} conserva solo las capturas indispensables: entrada manual, entrada de voz experimental, resultado P1 y registro de la decisión humana.

Para facilitar una evaluación exploratoria con participantes externos, el equipo preparó
también un acceso de prueba mediante un cuaderno en Google Colab. Este acceso no se
plantea como despliegue productivo ni como integración con sistemas 112, sino como
entorno reproducible de demostración: permite lanzar el prototipo, ejecutar casos
sintéticos y recoger respuestas mediante un cuestionario estructurado. La finalidad es
obtener valoración preliminar de utilidad, comprensibilidad, confianza, trazabilidad y
claridad de la supervisión humana sin exponer datos reales ni depender de una sala de
operaciones.

La integración generativa se realiza mediante Ollama con un modelo Llama 3.1 de 8.000 millones de parámetros, cuantizado para ejecución local. El LLM no decide la prioridad: recibe la
recomendación ya calculada por RuleFit-lite y la convierte en una explicación
estructurada. Si Ollama no responde, el prototipo conserva el modo degradado
determinista descrito en la sección siguiente.

Como comprobación técnica, se ha ejecutado una prueba funcional P1--P4 contra el
servicio web activo. Los resultados se conservan en los reportes reproducibles de evaluación
integrada del repositorio. La prueba verifica cuatro incidentes representativos, uno por
prioridad, con Capa 2 RuleFit activa y Capa 3 no degradada.

Adicionalmente, tras integrar la Capa 1, se ha regenerado una prueba técnica reproducible
Capa 1--Capa 2--Capa 3. Los informes resultantes se conservan como evidencias de
evaluación en el repositorio. Su finalidad es comprobar que la salida real de Capa 1
alimenta correctamente el motor RuleFit, incluyendo casos P1 con atrapamiento,
inconsciencia, incendio, intoxicación, inundación y una incidencia vial ordinaria con
negaciones explícitas como ``no hay heridos ni vehículos atrapados''.

\subsection{Interfaz de voz e inferencia inteligente}

Como demostración controlada, la interfaz incorpora un módulo experimental de entrada
por voz. El usuario puede grabar audio desde el micrófono o cargar un archivo de audio
de prueba; el reconocimiento automático del habla (ASR) se ejecuta localmente mediante
\texttt{faster-whisper} \cite{FasterWhisper2026}. La transcripción
resultante alimenta un parser estructurado que intenta extraer automáticamente título,
descripción, categoría, localidad y provincia. Cuando Ollama está disponible, el parser
emplea el modelo local \texttt{llama3.1:8b-instruct-q4\_K\_M} con temperatura 0.0; si el
modelo no responde, el sistema utiliza heurísticas locales basadas en palabras clave y
patrones de localidad o provincia.

El resultado de esta etapa no se envía directamente como decisión, sino que inicializa
reactivamente el estado de sesión de Streamlit mediante el patrón \textit{key-only}. El
operador conserva la posibilidad de revisar y corregir los campos antes de lanzar la
predicción. Esta funcionalidad pertenece al bloque de interfaz e integración del
prototipo y refuerza la orientación práctica del demostrador. No se ha evaluado con
llamadas reales, ruido de sala ni un corpus representativo de locutores, por lo que no
se presenta como un sistema de transcripción validado.

\begin{table}[ht]
\centering
\caption{Prueba funcional integrada del prototipo}
\label{tab:8-smoke-prototipo}
\small
\begin{tabular}{|l|c|c|r|c|}
\hline
\textbf{Caso} & \textbf{Esperada} & \textbf{Obtenida} & \textbf{Prob. dominante} & \textbf{LLM} \\ \hline
Incendio con atrapados & P1 & P1 & 97.35\% & Ollama \\ \hline
Fuga de gas sin heridos & P2 & P2 & 92.21\% & Ollama \\ \hline
Árbol en calzada & P3 & P3 & 93.71\% & Ollama \\ \hline
Consulta vecinal sin riesgo & P4 & P4 & 98.54\% & Ollama \\ \hline
\end{tabular}
\end{table}

La prueba supera los cuatro casos esperados y observa una latencia media de
\(4{,}898\) s por inferencia integrada, condicionada por el uso local del LLM. Esta
cifra no se interpreta como rendimiento operativo final, sino como primera evidencia de
integración funcional del prototipo completo. En el contexto de investigación académica
y ejecución sobre hardware doméstico con recursos compartidos, la medida sirve para
caracterizar el flujo completo, pero supera el objetivo interno de diseño. Una evolución
posterior debería estudiar hardware dedicado, una configuración específica del servidor
LLM o el uso de la respuesta alternativa determinista.

La prueba reproducible Capa 1--Capa 2--Capa 3 alcanza 5/5 casos correctos, con
modelo \texttt{RULEFIT} disponible, precisión micro de señales 0.733333, recall micro
1.000000 y F1 micro 0.846154 para Capa 1. Las latencias medias observadas son
4.58844 ms en Capa 1, 52.16898 ms en Capa 2 y 0.30748 ms en la explicación degradada
de Capa 3. Esta comprobación no sustituye a la prueba con el LLM local, pero deja evidencia
automática de que, tras la integración con Ancor, el sistema completo mantiene la
coherencia del flujo de datos.

\section{Contratos y trazabilidad}

El prototipo utiliza contratos versionados para separar responsabilidades entre capas. La salida de Capa 2 debe incluir, como mínimo, identificador del incidente, prioridad recomendada, probabilidades por clase, reglas o señales relevantes, versión del modelo y advertencias asociadas. La Capa 3 consume esa información y debe conservar la correspondencia entre explicación y evidencia.

La trazabilidad se apoya en evidencias reproducibles: conjunto de datos auditado, etiquetas débiles, particiones, modelos entrenados, métricas, errores P1, muestra preparada para revisión interna y síntesis final de evaluación. Este diseño permite reconstruir qué datos se usaron, qué modelo produjo la recomendación y qué explicación se generó.

\section{Modo degradado}

El prototipo contempla un modo degradado para la explicación. Si el LLM local no está disponible o no cumple el tiempo esperado, el sistema puede generar una explicación determinista basada en reglas, probabilidades y evidencias de la Capa 2. Este modo evita que la ausencia del componente generativo bloquee la recomendación o fuerce dependencia de servicios externos.

La respuesta alternativa determinista también ha permitido evaluar fidelidad explicativa de forma controlada sobre la muestra preparada de 119 casos, como se documenta en el Capítulo~\ref{cap:9}. En las pruebas reproducibles Capa 1--Capa 2--Capa 3, este modo ofrece respuestas submilisegundo y actúa como mecanismo de contingencia cuando el LLM local no está disponible o no cumple el tiempo esperado.

\section{Relación con la evaluación}

El prototipo no se evalúa por apariencia visual, sino por su capacidad para producir evidencias verificables. La evaluación experimental utiliza los mismos procedimientos, variables de entrada, particiones y modelos declarados por la arquitectura. Por eso los resultados del Capítulo~\ref{cap:9} pueden enlazarse con el diseño del Capítulo~\ref{cap:6}, los datos del Capítulo~\ref{cap:7} y los anexos de evaluación.

Un índice maestro estructurado actúa como registro de trazabilidad de esta relación entre prototipo, resultados y documentación. En este archivo de control se indexan las huellas de los conjuntos de datos de prueba, los pesos del modelo y las métricas calculadas. Esta estructura facilita la comprobación reproducible de los resultados declarados en la memoria.

\section{Repositorio del código fuente}

El código fuente completo, los scripts de datos, entrenamiento y evaluación, los contratos Pydantic y los artefactos versionados están disponibles en el repositorio público del proyecto junto a un notebook que permite la ejecución del prototipo en Google Colab. La URL de ambos son las siguientes:

\begin{center}
\url{https://github.com/jalbfil/tfe_muia_unir_188}
\end{center}


\begin{center}
\url{https://colab.research.google.com/drive/1XKVhjJCWDXH9bjI0Qsj4JCFUa5cwp7_J?usp=sharing}
\end{center}

\section{Conclusiones del capítulo}

El prototipo software desarrollado materializa la propuesta del TFM en una cadena ejecutable y auditable. Su diseño mantiene separadas la extracción, la priorización, la explicación y la validación humana, evitando cualquier tipo de fuga de información o leakage. Gracias a este enfoque, resulta factible comparar diversos modelos interpretables y dejar evidencias suficientes para defender los resultados de forma reproducible.

Con este prototipo, la evaluación del Capítulo~\ref{cap:9} no se apoya solo en una descripción teórica, sino en una implementación concreta de la arquitectura. Las métricas expuestas proceden de su ejecución sobre la partición de prueba. Los procedimientos y evidencias conservados permiten reproducir esos resultados en el entorno documentado, sin atribuirles validez operativa externa.
